\documentclass[12pt]{article}
\usepackage{tabularx}
\usepackage{amsmath} 
\usepackage{graphicx}
\usepackage[margin=1in]{geometry}
\usepackage{cite}
\usepackage[final]{hyperref}
\hypersetup{
	colorlinks=true,      
	linkcolor=blue,       
	citecolor=blue,       
	filecolor=magenta,    
	urlcolor=blue         
}
\usepackage{blindtext}
\usepackage{amssymb}
\usepackage{esint}
\usepackage{tikz}
\usepackage{pgfplots}
\usepackage{gensymb}
\usetikzlibrary{datavisualization}
\usetikzlibrary{datavisualization.formats.functions}
\newcommand{\Prho}{.8}%
\newcommand{\Ptheta}{55}%
\newcommand{\Pphi}{60}%
\usepackage{tikz-3dplot}
\newcommand{\uvec}[1]{\boldsymbol{\hat{\textbf{#1}}}}
%++++++++++++++++++++++++++++++++++++++++
\usepackage{empheq}

\newcommand*\widefbox[1]{\fbox{\hspace{2em}#1\hspace{2em}}}

\begin{document}

\usetikzlibrary{scopes}
\def\iangle{35} % Angle of the inclined plane

\def\down{-90}
\def\arcr{0.5cm} % Radius of the arc used to indicate angles


\title{ECE 3030 - HW 2}
\author{Joshua Tensuan - jvt24}
\date{September 9, 2022}
\maketitle


\begin{enumerate}
	\item [2.1]
	\begin{enumerate}
		\item [a.]
		Gauss' Law in integral form is given by,
		\[\oiint \epsilon_0 \vec E \cdot d\vec A = Q_{\mathrm{enc}}\]
		We can utilize cylindrical coordinates along the axis of the coaxial cylindrical surfaces.
		By the symmetries of the cylinders (reference lecture and previous homework for strict argument), we know that there cannot be any $\varphi, z$-dependence and no components in the $\varphi$ and $z$ direction.
		We can separate the space into three relevant regions: $r<a$, $a<r<b$, and $r>b$.
		Analyzing the first region, we know that $Q_{\mathrm{enc}}$ is 0 in this region for any concentric Gaussian cylinder formed with $r<a$.
		Thus, $\vec E = \vec 0$ in this region.
		Analyzing the second region, we can simplify the flux integral by the symmetries of the problem to,
		\[\oiint \epsilon_0\vec E \cdot d\vec A = \epsilon_0E_r\cdot A\]
		Inputting this into Gauss' Law in Integral Form,
		\[\epsilon_0E_r\cdot A = \sigma_{s,a}A_a\to \epsilon_0E_r\cdot2\pi rl = \sigma_{s,a}2\pi a l\]
		\[E_r = \frac{\sigma_{s,a}}{\epsilon_0}\frac{a}{r}\]
		Thus, in the region, $\vec E$ is given by,
		\[\vec E = \frac{\sigma_{s,a}}{\epsilon_0}\frac{a}{r}\hat r\]
		Analyzing the last region, we can again simplify the flux integral by the symmetries of the problem to,
		\[\oiint \epsilon_0 \vec E \cdot d\vec A = \epsilon_0 E_r\cdot A\]
		Inputting this into Gauss' Law in Integral Form,
		\[\epsilon_0E_r\cdot A = \sigma_{s,a}A_a + \sigma_{s,b}A_b\to \epsilon_0E_r\cdot 2\pi rl = \sigma_{s,a}\cdot2\pi al+\sigma_{s,b}\cdot 2 \pi bl\]
		\[E_r = \frac{\sigma_{s,a}}{\epsilon_0}\frac{a}{r}+\frac{\sigma_{s,b}}{\epsilon_0}\frac{b}{r}\]
		Thus, in the region, $\vec E$ is given by,
		\[\vec E = \left(\frac{\sigma_{s,a}}{\epsilon_0}\frac{a}{r}+\frac{\sigma_{s,b}}{\epsilon_0}\frac{b}{r}\right) \hat r=\frac{1}{\epsilon_0}\cdot\frac{1}{r}\left(\sigma_{s,a}a + \sigma_{s,b}b\right)\hat r\]
		\[\boxed{\vec E = \begin{cases}
			\vec 0 & r < a \\ 
			\frac{\sigma_{s,a}}{\epsilon_0}\frac{a}{r}\hat r & a < r < b \\ 
			\frac{1}{\epsilon_0}\cdot \frac{1}{r}(\sigma_{s,a}a + \sigma_{s,b}b)\hat r & r > b
		\end{cases}}\]

		\item [b.]
		If $\vec E$ vanishes when $r>b$, then $\vec E = \vec 0$ in the last region considered in the last part.
		\[\frac{1}{\epsilon_0}\cdot\frac{1}{r}\left(\sigma_{s,a}a + \sigma_{s,b}b\right)\hat r= \vec 0\]
		\[\to \sigma_{s,a}a + \sigma_{s,b}b = 0 \to b = \frac{\sigma_{s,a}}{\sigma_{s,b}}a\]
		Thus, in order for $\vec E$ to vanish when $r>b$, then $a$ and $b$ must be related by,
		\[\boxed{b = \frac{\sigma_{s,a}}{\sigma_{s,b}}a}\]
	\end{enumerate}

	\item [2.2]
	\begin{enumerate}
		\item [a.]
		The $z$-component of the electric field as a position $x$ along the $x$-axis (where $z=0$) must be 0.
		This is due to the symmetries of the problem.
		We can show this utilizing a proof by contradiction.
		Without loss of generality (works regardless of sign), assume there is an electric field along the $+z$-direction at any point along the $x$-axis.
		Then, flip the entire system (included the associated electric fields) about the $x$-axis.
		Because nothing fundamental was changed about the system, there will now be an electric field along the $-z$-direction where there existed an electric field along the $+z$-direction in the original system.
		However, we notice that the original system is exactly the same as the flipped system.
		However, there is supposedly an electric field in both the $-z$ and $+z$-direction, creating a contradiction.
		Thus, there cannot be an electric field in the $z$-direction at any point along the $x$-axis.
		\[\boxed{\vec E(x, z=0) = \vec 0}\]

		\item [b.]
		From Coulomb's Law, we know that for an infinitesimal point charge,
		\[d\vec E = \frac{1}{4\pi\epsilon_0}\frac{dq}{r^2}\hat r\]
		where $r$ is the displacement vector between the infinitesimal point charge and the point we evaluate the electric field at.
		We can express $r$ as, $r^2=x^2+z^2$.
		Additionally, from part a, we know that the $z$-component of the electric field is 0.
		Thus, we only have to solve for the $x$-component of the electric field.
		Taking the projection of the infinitesimal electric field along the $x$-axis.
		\[dE_x = d\vec E \cdot \hat x = \frac{1}{4\pi\epsilon_0}\frac{dq}{x^2+z^2}\left(\frac {x}{\sqrt{x^2+z^2}}\right)\]
		From the prompt, we know that $dq = \lambda \: dz$.
		Substituting this in,
		\[dE_x=\frac{1}{4\pi\epsilon_0}\frac{\lambda \: dz}{x^2+z^2}\left(\frac{x}{\sqrt{x^2+z^2}}\right)\]
		Integrating across the line charge,
		\[E_x = \int_{-L/2}^{L/2}\frac{\lambda}{4\pi\epsilon_0}\frac{x}{(x^2+z^2)^{3/2}}\: dz\]
		Factoring out constant terms and using the symmetry of the integral, we can simplify the expression to,
		\[E_x = \frac{\lambda}{2\pi\epsilon_0}\int_0^{L/2}\frac{x}{(x^2+z^2)^{3/2}}\: dz\]
		Performing the substitution of $z=x\tan(u)$, $dz = x\sec^2(u)\: du$,
		\[E_x = \frac{\lambda}{2\pi\epsilon_0}\int_0^{\arctan(L/2x)}\frac{x}{(x^2+x^2\tan^2(u))^{3/2}}\cdot x\sec^2(u)\: du\]
		Through some algebra and use of trigonometric identities,,
		\[E_x = \frac{\lambda}{2\pi\epsilon_0}\int_0^{\arctan(L/2x)}\frac{ x^2\sec^2(u)}{(x^2\sec^2(u))^{3/2}}\: du\]
		\[E_x = \frac{\lambda}{2\pi\epsilon_0}\int_0^{\arctan(L/2x)}\frac{\cos(u)}{x}\: du\]
		\[E_x = \frac{\lambda}{2\pi\epsilon_0}\left[\frac{1}{x}\sin(u)\right]_{u=0}^{u=\arctan(L/2x)}=\frac{\lambda}{2\pi\epsilon_0}\left(\frac{L}{x\sqrt{L^2+4x^2}}\right)\]
		\[\boxed{E_x=\frac{\lambda}{2\pi\epsilon_0}\left(\frac{L}{x\sqrt{L^2+4x^2}}\right)}\]

		\item [c.]
		We are interested in taking the limit that $L\ggg x$.
		In this case, we can approximate $\sqrt{L^2+4x^2}\approx \sqrt{L^2} = L$.
		Substituting this into $E_x$,
		\[E_x\approx \frac{\lambda}{2\pi \epsilon_0}\frac{L}{xL}=\frac{\lambda}{2\pi\epsilon_0}\frac{1}{x}\]
		\[\boxed{E_x \approx \frac{\lambda}{2\pi\epsilon_0}\frac{1}{x}}\]

		\item [d.]
		Our answer here is equivalent to the answer in problem 1.5(c) if we projected the wire onto a two-dimensional plane.
		Here, the $z$-axis stays the same, $x$-axis points in radial direction.
		This is because in the limit $L\ggg x$, the line of charge is approximately the same as the two-dimensional projection of an infinite wire of charge.

		\item [e.]
		We are now interested in the limit that $x \ggg L$.
		In this case, we can approximate $\sqrt{L^2+4x^2}\approx \sqrt{4x^2}=2x$.
		Substituting this into $E_x$,
		\[E_x \approx \frac{\lambda}{2\pi \epsilon_0}\frac{L}{x\cdot 2x}=\frac{\lambda}{4\pi \epsilon_0}\frac{L}{ x^2}\]
		\[\boxed{E_x \approx \frac{\lambda}{4\pi \epsilon_0}\frac{L}{ x^2}}\]
		This expression is reminiscent to that of the electric field of a point charge ($x$ represents projection of radial direction here).
		This makes intuitive sense as far away from the line of charge, the configuration should 'look' like a point charge.
		Here, we can make a direct comparison by setting $\lambda L = Q$ of the configuration.
		\[\boxed{\text{Expression looks like that of a point charge.}}\]

	\end{enumerate}

	\item [2.3]
	\begin{enumerate}
		\item [a.]
		\begin{enumerate}
			\item [i.]
			Here, we can form a concentric Gaussian spherical shell with radius $0<r<b$.
			The resultant electric field from Gauss' Law is then simply just the electric field from a point charge.
			Thus,
			\[\boxed{\vec E(\vec r)=\frac{1}{4\pi \epsilon_0}\frac{q}{\|\vec r\|^2}\hat r, \; 0<r<b}\]

			\item [ii.]
			We know that in any conducting material, the electric field must be 0.
			Otherwise, there will be an infinite current density by Ohm's Law which is non-physical.
			Thus,
			\[\boxed{\vec E(\vec r)=\vec 0, \; b<r<c}\]

			\item [iii.]
			We can again draw a concentric Gaussian spherical shell with radius $c<r<\infty$.
			Due to the conservation of charge, we know that the enclosed charge must still be $+q$.
			The resultant electric field from Gauss' Law is then simply just the electric field from a point charge.
			Thus,
			\[\boxed{\vec E(\vec r)=\frac{1}{4\pi \epsilon_0}\frac{q}{\|\vec r\|^2}\hat r, \; c<r<\infty}\]

		\end{enumerate}

		\item [b.]
		From class, we know that we must follow the boundary condition of,
		\[\epsilon_0(E_2-E_1)=\sigma\]
		This is trivially derived from using Gauss' Law on a small area of the surface.
		First applying this on the inner surface.
		We can use our previous expression from part a to find that just inside the inner surface, the electric field magnitude is given by,
		\[E_1 = \frac{1}{4\pi \epsilon_0}\frac{q}{b^2}\]
		We know that inside the conductor, $E_2=0$.
		Substituting this in,
		\[\epsilon_0\left(0-\frac{1}{4\pi \epsilon_0}\frac{q}{b^2}\right)=\sigma_i\]
		where $\sigma_i$ is the inner surface charge density.
		Solving for $\sigma_i$,
		\[\boxed{\sigma_i = -\frac{1}{4\pi}\frac{q}{b^2}}\]
		Now applying this on the outer surface.
		We can use our previous expression from part a to find that just outside the inner surface, the electric field magnitude is given by,
		\[E_2 = \frac{1}{4\pi\epsilon_0}\frac{q}{c^2}\]
		We know that inside the conductor, $E_1=0$.
		Substituting this in,
		\[\epsilon_0\left(\frac{1}{4\pi\epsilon_0}\frac{q}{c^2}-0\right)=\sigma_o\]
		where $\sigma_o$ is the outer surface charge density.
		Solving for $\sigma_o$,
		\[\boxed{\sigma_o = \frac{1}{4\pi}\frac{q}{c^2}}\]
		The total charge is then given by,
		\[Q = \sigma_oA_o+\sigma_iA_i\]
		where $A_o$ and $A_i$ are the outer and inner surface areas of the shell respectively and $Q$ is the total charge.
		Substituting values in,
		\[Q=\frac{1}{4\pi}\frac{q}{c^2}\cdot 4\pi c^2 - \frac{1}{4\pi}\frac{q}{b^2}\cdot 4\pi b^2=0\]
		\[\boxed{Q=0}\]

		\item [c.]
		\begin{enumerate}
			\item [i.]
			We can solve for Laplace's equation in the region around the point charge.
			Because there are no charges here,
			\[\nabla^2\phi = 0\]
			Using the spherical symmetries of the problem (rotationally invariant), we can find that,
			\[\phi(r) = \frac{A}{r}+ B\]

			\item [ii.]
			Because there are no charges within the shell, (can ignore surface charges)
			\[\nabla^2\phi = 0\]
			Using the spherical symmetries of the problem (rotationally invariant), we can find that,
			\[\phi(r) = \frac{C}{r}+ D\]
			Because we know that there is no electric field in the shell, we know that $\phi$ must be constant, i.e $\phi(r)=D$.

			\item [iii.]
			Because there are no charges outside of the shell, 
			\[\nabla^2\phi = 0\]
			Using the spherical symmetries of the problem (rotationally invariant), we can find that,
			\[\phi(r) = \frac{E}{r}+F\]
			We know that at $\phi(\infty)=0$.
			Thus, $F=0$.

		\end{enumerate}
		We then end up with three equations with four unknowns.
		Using the fact that the potential must be continuous on the outer surface of the shell,
		\[D = \phi(c)\to D = \frac{E}{c}\]
		Using the fact that the potential must be continuous on the inner surface of the shell,
		\[\frac{A}{b}+B=D=\frac{E}{c}\]
		Now applying boundary conditions of discontinuities in the electric field.
		We can first consider the inner shell.
		We know that $\epsilon_0(E_2-E_1)=\sigma_i$.
		\[\begin{cases}
			E_1 = -\partial \phi_1/\partial r|_{r=b} = A/b^2 \\
			E_2 = -\partial \phi_2/\partial r = 0
		\end{cases}\]
		where $\phi_i$ refers to the potential in the $i$-th region.
		Then,
		\[\epsilon_0(E_2-E_1)=\sigma_i \to -\epsilon_0\frac{A}{b^2}=-\frac{1}{4\pi}\frac{q}{b^2}\]
		\[A=\frac{q}{4\pi \epsilon_0}\]
		Similarly, considering the outer shell, we know that $\epsilon_0(E_2-E_1)=\sigma_o$.
		\[\begin{cases}
			E_1 = -\partial \phi_2/\partial r= 0 \\	
			E_2 = -\partial \phi_3/\partial r|_{r=c} = E/c^2
		\end{cases}\]
		Then,
		\[\epsilon_0(E_2-E_1)=\sigma_o \to \epsilon_0\frac{E}{c^2}=\frac{1}{4\pi}\frac{q}{c^2}\]
		\[E = \frac{q}{4\pi\epsilon_0}\]
		Using our continuous potential condition,
		\[\frac{q}{4\pi\epsilon_0}\cdot \frac{1}{b} + B = \frac{q}{4\pi\epsilon_0}\cdot \frac{1}{c}\]
		\[B = \frac{q}{4\pi\epsilon_0}\left(\frac{1}{c}-\frac{1}{b}\right)\]
		Having solved for all the unknowns, we can substitute this back into the equations to yield,
		\begin{enumerate}
			\item [i.]
			\[\boxed{\phi(r)= \frac{q}{4\pi\epsilon_0}\frac{1}{r}+\frac{q}{4\pi\epsilon_0}\left(\frac{1}{c}-\frac{1}{b}\right), \; 0<r\leq b}\]

			\item [ii.]
			\[\boxed{\phi(r)=\frac{q}{4\pi\epsilon_0}\frac{1}{c}, \; b\leq r\leq c}\]

			\item [iii.]
			\[\boxed{\phi(r)=\frac{q}{4\pi\epsilon_0}\frac{1}{r}, \; c\leq r \leq \infty}\]
		\end{enumerate}

		\item [d.]
		Within the conductor, the electric field must be 0.
		Otherwise, there will be an infinite current density by Ohm's Law which is non-physical.
		Inside the conducting shell, we can utilize Gauss' Law, and we would get the same electric field as gotten in part a.
		Outside the conducting shell, we know that $\phi(c)=\phi(\infty)=0$.
		In order to satisfy Laplace's Equation, there cannot be local minima or maxima in the potential in the region outside the conducting shell.
		Thus, $\phi$ must be $0$ at all points outside the conducting shell.
		By extension, $\vec E$ must be 0 at all points outside the conducting shell.
		\begin{enumerate}
			\item [i.]
			\[\boxed{\vec E(\vec r)=\frac{1}{4\pi\epsilon_0}\frac{q}{\|\vec r\|^2}\hat r, \; 0 < r < b}\]

			\item [ii.]
			\[\boxed{\vec E(\vec r)=\vec 0, \; b < r < c}\]

			\item [iii.]
			\[\boxed{\vec E(\vec r)=\vec 0, \; c < r < \infty}\]
		\end{enumerate}

		\item [e.]
		From class, we know that the system must follow the boundary condition of, 
		\[\epsilon_0(E_2-E_1)=\sigma\]
		This is trivially derived from using Gauss' Law on a small area of the surface.
		First applying this on the inner surface.
		We can use our previous expression from part a to find that just inside the inner surface, the electric field magnitude is given by,
		\[E_1 = \frac{1}{4\pi \epsilon_0}\frac{q}{b^2}\]
		We also know that $E_2=0$.
		Substituting this in,
		\[\epsilon_0\left(0 - \frac{1}{4\pi\epsilon_0}\frac{q}{b^2}\right)=\sigma_i\]
		where $\sigma_i$ is the inner surface charge density.
		Solving for $\sigma_i$,
		\[\boxed{\sigma_i = -\frac{1}{4\pi}\frac{q}{b^2}}\]
		Now applying this on the outer surface, we know $E_1=E_2=0$.
		Thus, $\sigma_o=0$.
		\[\boxed{\sigma_o =0}\]
		The total charge is then given by,
		\[Q=\sigma_oA_o+\sigma_iA_i\]
		where $A_o$ and $A_i$ are the outer and inner surface areas of the shell respectively and $Q$ is the total charge.
		Substituting values in,
		\[Q=0-\frac{1}{4\pi}\frac{q}{b^2}\cdot 4\pi b^2=-q\]
		\[\boxed{Q=-q}\]

		\item [f.]
		\begin{enumerate}
			\item [i.]
			We can solve for Laplace's equation in the region around the point charge.
			Because there are no charges here,
			\[\nabla^2\phi = 0\]
			Using the spherical symmetries of the problem (rotationally invariant), we can find that,
			\[\phi(r)=\frac{A}{r}+B\]

			\item [ii.]
			Because there are no charges within the shell, (can ignore surface charges)
			\[\nabla^2\phi = 0\]
			Using the spherical symmetries of the problem (rotationally invariant), we can find that,
			\[\phi(r)=\frac{C}{r}+D\]
			Because we know that there is no electric field in the shell, we know that $\phi$ must be constant, i.e $\phi(r)=D$.
			
			\item [iii.]
			Because there are no charges outside of the shell,
			\[\nabla^2 \phi = 0\]
			Using the spherical symmetries of the problem (rotationally invariant), we can find that,
			\[\phi(r)=\frac{E}{r}+F\]
			We know that at $\phi(\infty)=0$. 
			Thus, $F=0$.
			Additionally, we know that because there is no electric field outside the shell, we know that $\phi$ must be constant, i.e $\phi(r)=0$.

		\end{enumerate}

		We can use the boundary condition that potential must be continuous.
		Analyzing continuity of the potential at the outer surface,
		\[\phi_2(c)=\phi_3(c) \to D=0\]
		where $\phi_i$ represents the potential in the $i$-th region.
		Thus, $\phi_2 = \phi_3 = 0$.
		Analyzing continuity of the potential at the inner surface,
		\[\phi_1(b)=\phi_2(b)\to \frac{A}{b}+B = 0\]
		Now applying boundary conditions of the discontinuities in the electric field.
		We can consider the inner shell.
		We know that $\epsilon_0(E_2-E_1)=\sigma_i$.
		\[\begin{cases}
			E_1 = -\partial \phi_1/\partial _r |_{r=b}=A/b^2 \\ 
			E_2 = -\partial\phi_2/\partial_r = 0	
		\end{cases}\]
		where $\phi_i$ again refers to the potential in the $i$-th region.
		Then,
		\[\epsilon_0(E_2-E_1)=\sigma_i \to -\epsilon_0\frac{A}{b^2}=-\frac{1}{4\pi}\frac{q}{b^2}\]
		\[A =\frac{q}{4\pi\epsilon_0}\]
		Substituting this into the continuity condition,
		\[\frac{q}{4\pi\epsilon_0}\frac{1}{b}+B=0\to B = -\frac{q}{4\pi\epsilon_0}\frac{1}{b}\]
		Substituting this back into the potentials,
		\begin{enumerate}
			\item [i.]
			\[\boxed{\phi(r)=\frac{q}{4\pi\epsilon_0}\left(\frac{1}{r}-\frac{1}{b}\right)}\]

			\item [ii.]
			\[\boxed{\phi(r)=0}\]

			\item [iii.]
			\[\boxed{\phi(r)=0}\]
		\end{enumerate}

		\item [g.]
		We notice that in both parts, we have the same charge induced on the inner shell.
		However, in part b, we have the same magnitude but opposite charge induced on outer shell while in part e, we have no charge induced.
		Grounding the shell allowed an outlet for negative charge to flow in (conversely, positive charge to flow out).
		Thus, in part b, charge conservation was conserved, while in part e, we notice that there was a non-zero net charge on the conductor.
		Grounding the shell allowed charge to flow into the system.

		\item [h.]
		We notice that inside the conductor, the potential is shifted by a constant factor in part f (lower potential).
		Additionally, in part f, the conducting shell and outside the conducting shell has a potential of 0.
		In other words, grounding the conducting shell forced the potential of the conducting shell to be the same as the potential at infinity.
		In other words, grounding the conducting shell allowed potential energy to leave the system.
		By allowing negative charge to flow into the system, the overall configuration lost electric potential.
		
	\end{enumerate}

	\item [4.]
	\begin{enumerate}
		\item [a.]
		For the region $z>0$, we can use the method of images.
		Because the conductor is at an equipotential, we know that for all points of $z=0$, the potential will be the same.
		We can let this potential be $0$.
		This gives us Laplace's Equation with the given boundary condition.
		If we place an image charge of $-q$ at $z=-d/2$, then Laplace's equation for the region $z>0$ and the given boundary condition $\phi(z=0)=0$ will be satisfied.
		By the uniqueness theorem, we know that this must be the solution for $z>0$.

		We can now evaluate the potential for the region $z>0$ (equivalent to derivation of potential from dipole done in class).
		Let $r_+$ be the distance between $+q$ and the point we evaluate the potential at and $r_-$ be the distance between the image charge $-q$ and the point we evaluate the potential at.
		Then we can take a superposition of the two potentials to yield,
		\[\phi(\vec r)=\frac{q}{4\pi\epsilon_0 r_+}-\frac{q}{4\pi\epsilon_0 r_-}\]
		We can use the approximation that $r_+\approx r-(d/2)\cos(\theta)$ and $r_-\approx r+(d/2)\cos(\theta)$.
		Then,
		\[\phi(\vec r)\approx\frac{q}{4\pi\epsilon_0\left(r-\frac{d}{2}\cos(\theta)\right)}-\frac{q}{4\pi\epsilon_0\left(r+\frac{d}{2}\cos(\theta)\right)}\]
		Through some algebraic manipulation,
		\[\phi(\vec r)\approx \frac{q\left(r+\frac{d}{2}\cos(\theta)\right)-q\left(r-\frac{d}{2}\cos(\theta)\right)}{4\pi\epsilon_0\left(r-\frac{d}{2}\cos(\theta)\right)\left(r+\frac{d}{2}\cos(\theta)\right)}\]
		\[\phi(\vec r)\approx \frac{qd\cos(\theta)}{4\pi\epsilon_0\left(r^2-\frac{d^2}{4}\cos^2(\theta)\right)}\]
		Using the approximation that for $r\gg d$, $r^2-\frac{d^2}{4}\cos^2(\theta)\approx r^2$,
		\[\phi(\vec r)\approx\frac{qd}{4\pi\epsilon_0 r^2}\cos(\theta), \; z>0\]

		Now evaluating Laplace's equation for the region $z<0$.
		We know that the spherical solution is given by $\phi(r)=\frac{A}{r}+B$.
		Since the whole perfect conducting slab must be an equipotential, there cannot be any variance in $\phi$ within the slab.
		Necessarily, this means that $A$ must be 0.
		Additionally, we can use the boundary condition that $\phi(z=-d/4)=0$ (conductor slab is equipotential) to find that $B$ must also be 0.
		Because this satisfies Laplace's equation with the given boundary conditions, we know that by the uniqueness theorem, this must be the solution.
		Thus,
		\[\phi(\vec r) =0, \; z< 0\]
		\[\boxed{\phi(\vec r) = \begin{cases}
			\sim \frac{qd}{4\pi\epsilon_0 r^2}\cos(\theta), & z> 0 \\ 
			0, & z<0
		\end{cases}}\]

		\item [b.]
		We know that $\vec E = -\nabla \phi$.
		Trivially, we know that for $z<0$, $\vec E=-\nabla(0) = 0$.
		Solving for $\vec E$ when $z>0$,
		\[\vec E = -\nabla \phi \approx - \nabla\left(\frac{qd}{4\pi\epsilon_0r^2}\cos(\theta)\right)\]
		For a function $f(r,\theta)$ with only $r$ and $\theta$ dependence in spherical coordinates,
		\[\nabla f = \frac{\partial f}{\partial r}\hat r + \frac{1}{r}\frac{\partial f}{\partial \theta}\hat \theta\]
		Applying this to $\phi$,
		\[-\nabla \phi = \frac{qd}{2\pi\epsilon_0r^3}\cos(\theta)\hat r + \frac{qd}{4\pi\epsilon_0r^3}\sin(\theta)\hat \theta\]
		\[\boxed{\vec E = \frac{qd}{2\pi\epsilon_0r^3}\cos(\theta)\hat r + \frac{qd}{4\pi\epsilon_0r^3}\sin(\theta)\hat \theta}\]

		\item [c.]
		From class, we have the boundary condition that,
		\[\epsilon_0 ( E_{2,\perp} - E_{1,\perp})=\sigma\]
		At $z=0$, $\theta=\pi/2$ for $+x$ and $\theta = -\pi/2$ for $-x$.
		In these conditions (on $xy$-plane), we cannot make the same approximations as done previously.
		Thus, we must calculate the electric fields from scratch.
		We can utilize Coulomb's Law to find the electric fields from each of the charges.
		Through the principle of superposition, we can sum the electric fields to find the total electric field.
		We know that in the static case, there cannot be any electric fields parallel to the conductor boundary.
		Thus, the only electric field will be perpendicular to the conductor, and we only have to analyze $E_z$.
		From the top charge,
		\[E_{z,1} = -\frac{1}{4\pi\epsilon_0}\frac{q}{(r^2+(d/2)^2}\cdot \frac{d/2}{\sqrt{r^2+(d/2)^2}}=-\frac{qd}{8\pi\epsilon_0\left(r^2+(d/2)^2\right)^{3/2}}\]
		Similarly for the image charge,
		\[E_{z,2}=\frac{1}{4\pi\epsilon_0}\frac{-q}{(r^2+(d/2)^2}\cdot \frac{d/2}{\sqrt{r^2+(d/2)^2}}=-\frac{qd}{8\pi\epsilon_0\left(r^2+(d/2)^2\right)^{3/2}}\]
		Thus, by the principle of superposition,
		\[E_z = E_{z,1}+E_{z,2}=-\frac{qd}{4\pi\epsilon_0(r^2+(d/2)^2)^{3/2}}\]
		
		We know that inside the  conductor $E_{1,\perp}=0$.
		Substituting this into the perpendicular electric field boundary condition,
		\[\sigma_u = -\epsilon_0\frac{qd}{4\pi\epsilon_0(r^2+(d/2)^2)^{3/2}}=-\frac{qd}{4\pi(r^2+(d/2)^2)^{3/2}}\]
		where $\sigma_u$ is the charge density on the upper surface on the metal strip.
		\[\boxed{\sigma_u = -\frac{qd}{4\pi(r^2+(d/2)^2)^{3/2}}}\]
		For the bottom surface, we know that $E_{2,\perp}=E_{1,\perp}=0$.
		Thus, $\sigma_l=0$ where $\sigma_l$ is the surface charge density on the lower surface.
		\[\boxed{\sigma_l = 0}\]

		\item [d.]
		To find the total charge on the upper surface of the metal strip, we just have to integrate the surface charge density over the entire area of the strip.
		\[Q_u= \iint_S \sigma_u \: dA\]
		where $Q_u$ denotes the total charge on the upper surface of the metal strip.
		Using spherical coordinates,
		\[Q_u = \int_0^{2\pi}\int_0^{\infty}-\frac{qd}{4\pi (r^2+(d/2)^2)^{3/2}}r \: dr \: d\theta\]
		Integrating over $r$,
		\[Q_u=\int_0^{2\pi}\left[\frac{qd}{4\pi(r^2+(d/2)^2)^{1/2}}\right]_0^{\infty}\: d\theta\]
		Enforcing the bounds,
		\[Q_u=\int_0^{2\pi}-\frac{qd}{4\pi(0^2+(d/2)^2)^{1/2}}\:d\theta\]
		Through algebraic manipulation,
		\[Q_u=\int_0^{2\pi}-\frac{q}{2\pi}\:d\theta=-q\]
		\[\boxed{Q_u=-q}\]






	\end{enumerate}
	
	
\end{enumerate}

		

\end{document}
